\emph{Mesin Turing} minangka mékanisme komputasi kang
diidéalisasi. Mesin iki nduwèni \emph{pita} tanpa wates ing
siji arah, kapérang dadi kothak-kothak. Saben kothak bisa ngemot
\emph{simbol} saka alfabèt kang wis ditetepaké. Mesin lumaku
kanthi ngobahaké \emph{sirah waca-tulis} ing sadawané pita.
Mesin uga ana ing salah siji saka \emph{kahanan} kang wis
ditetepaké. Pakaryan sirah waca-tulis ditentokaké déning
himpunan prentah. Saben prentah ditrapaké nalika mesin ana
ing kahanan tartamtu lan maca simbol tartamtu. Prentah iku
nemtokaké simbol kang ditulis ing kothak kang lagi diwaca,
apa sirah obah sakothak mangiwa, manengen, utawa tetep, lan
menyang kahanan endi mesin pindhah. \emph{Masukan} diwujudaké
minangka urutan simbol ing pita. Mesin diwiwiti ing kahanan
wiwitan, kanthi sirah waca-tulis maca kothak masukan paling
kiwa. Himpunan prentah nemtokaké urutan \emph{konfigurasi}
saben langkah, yaiku isi pita, papané sirah waca-tulis, lan
kahanan mesin. Yèn ing sawijining konfigurasi ora ana
prentah kanggo pasangan simbol kang lagi diwaca lan kahanan
kang lagi dienggo, mesin \emph{mandheg}. Isi pita dadi
\emph{asil}.

Wilangan gampang diwujudaké minangka urutan garis ing pita
mesin Turing. Fungsi $f\colon\Nat\to\Nat$ \emph{bisa
dikomputasi déning mesin Turing} yèn ana mesin kang, saben
diwiwiti nganggo representasi uner~$n$ minangka masukan,
pungkasané mandheg kanthi representasi uner~$f(n)$ ing pita
minangka asil. Akèh fungsi aritmetika kang wis lumrah bisa
dibuktèkaké bisa dikomputasi déning mesin Turing, sanajan
sawetara buktiné luwih ruwet. Akèh modèl komputasi liya wis
diusulaké. Fungsi-fungsi aritmetika kang bisa dikomputasi ing
modèl mau tansah bisa dikomputasi déning mesin Turing.
Iki dadi panyengkuyung \emph{Tèsis Church--Turing}, yaiku yèn
saben fungsi aritmetika kang bisa dikomputasi kanthi èfèktif
uga bisa dikomputasi déning mesin Turing.
